\documentclass[10pt,a4paper]{amsart}
\usepackage[margin=19mm]{geometry}
\usepackage{amsmath,amssymb,mathtools}
\usepackage[scr=rsfs,cal=euler]{mathalfa}
\usepackage{xfrac}
\usepackage[unicode,hidelinks,bookmarks=false]{hyperref}
\newcommand{\ie}{i.e.\ }
\newcommand{\cf}{cf.\ }
\newcommand{\resp}{respectively\ }
\newcommand{\Subsection}{\subsection}
\providecommand{\nobreakdash}{\nobreak\hbox{-}}
\setlength{\emergencystretch}{2em}
\multlinegap0pt
\usepackage{mathtools}
\usepackage{amssymb}
\usepackage{dsfont}
\newtheorem{proposition}{Proposition}
\newcommand{\C}{\mathds{C}}
\newcommand{\R}{\mathds{R}}
\newcommand{\Rp}{\R^{+}}
\renewcommand{\d}{\mathrm{d}}
\title{Cutoff for geodesic paths on hyperbolic manifolds}
\author{Charles Bordenave, Joffrey Mathien}
\date{Source: arXiv:2502.06325v2 (CC BY 4.0)}
\begin{document}
\maketitle

\noindent\emph{Source and licence.} Cutoff for geodesic paths on hyperbolic manifolds. Version arXiv:2502.06325v2 (CC BY 4.0), distributed by arXiv under the Creative Commons Attribution 4.0 International licence, http://creativecommons.org/licenses/by/4.0/, as recorded in the arXiv licence metadata for this exact version and shown on the abstract page https://arxiv.org/abs/2502.06325v2. Full licence notice: \texttt{licenses/arxiv-2502.06325-CC-BY-4.0.txt}.\par\smallskip

\noindent\textit{Editorial scope.} Counted unit: the article's proposition prop:SpAt (source0.tex lines 312--319). The article's complete original proof of it is delivered with it (source0.tex lines 409--539, one \texttt{proof} environment of 131 lines, which the article places away from the statement). No mathematics is written, repaired or re-derived: the statement and the proof are the article's own lines, in the article's own notation, and no retained line is substituted or re-flowed.  Retained uncounted as the source-local context the proof needs: lines 299--310.  Nothing was omitted from any retained span, so the package carries no omission record.  No retained line contains a citation, so the wrapper carries no bibliography and no bibliography entry.  No retained line contains a figure environment, an image-inclusion command or drawing code, so no image is delivered and the package declares no asset dependency.  Editorial formatting only, in addition: an AMS \texttt{amsart} preamble that restates the article's own declarations and macros used by the retained lines, namely the environments \texttt{proposition} on separate counters, together with the standard packages those lines need.  The retained environments print in this document as Proposition 1 in order of appearance, whereas the article prints the same items with the numbering of its own full document.  The complete native source of the article as distributed in the arXiv e-print for this exact version is bundled unchanged as \texttt{sources/173025/source0.tex}.
\begin{proposition}{}\label{prop:SpAt0}
	Let $\phi_k$ be an eigenfunction of the Laplace-Beltrami operator, with eigenvalue $\lambda_k$. Then, $\phi_k$ is also an eigenfunction of $A_t$, and there exists a function $\nu_t\colon \Rp \to \R$ such that \[A_t\phi_k = \nu_t(\lambda_k)\phi_k.\]
	If $\lambda = \frac{(d-1)^2}{4}+u^2$, $u \in \R \cup i \left[-\frac{d-1}{2}, \frac{d-1}{2}\right]$,
	we have 
$$
\nu_t(\lambda) =  c_d \int_{-1}^1 ( \cosh t  +  x \sinh t  )^{iu - \frac{d-1}{2}}  ( 1 - x^2)^{ \frac{d - 3}{2}} \d x,  
$$
where $c_d = \frac{\Gamma\left( \frac{d}{2}\right)}{\sqrt{\pi} \Gamma\left( \frac{d-1}{2}\right)}$. In particular, for $d=3$ 
$$
\nu_t(\lambda) = \frac{\sin (ut)}{u \sinh (t)}.
$$
\end{proposition}

\begin{proposition}{}\label{prop:SpAt}
	Let $\lambda = \frac{(d-1)^2}{4}+u^2$ with $u \in \R_+ \cup i \left[0, \frac{d-1}{2}\right]$. Let $t_0 > 0$. There exists a constant $C = C(d,t_0)$ such that for all $t \geq t_0$, if $u \in i \left[0, \frac{d-1}{2}\right]$, we have
	$$|\nu_t(\lambda)|  \leq C  \min \left( t, \frac{1}{\Im(u)} \right)  e^{ \Im (u)  t-\frac{d-1}{2}t},$$
	and  if $u \in \R_+$,
	$$
	|\nu_t(\lambda)|  \leq C  \min \left( t, \frac{1}{u^{\frac{d-1}{2}}} \right)  e^{-\frac{d-1}{2}t}.
	$$
\end{proposition}

\begin{proof}[Proof of Proposition \ref{prop:SpAt}]
	For $d=3$, the lemma is clear from the explicit expression given in Proposition \ref{prop:SpAt0}. For general $d \geq 2$, we start by proving that all $u \in \R_+ \cup i \left[0, \frac{d-1}{2}\right]$,
	\begin{equation}\label{eq:bdnut1}
		|\nu_t(\lambda)|  \leq C  \min \left( t, \frac{1}{\Im(u)} \right)  e^{ \Im (u)  t-\frac{d-1}{2}t}.
	\end{equation}
	Let $v = \Im(u) \geq 0$ and $a = \tanh t \in [\tanh t_0,1)$. By Proposition \ref{prop:SpAt0} (applied to $-u$), 
	\begin{align*}
		|  \nu_t (\lambda) | &= c_d (\cosh t )^{v - \frac{d-1}{2}} \left| \int_{-1}^{1} ( 1+ a x)^{-iu - \frac{d-1}{2}} ( 1- x^2)^{ \frac{d - 3}{2}} \d x \right| \\
		& \leq c_d (\cosh t )^{v - \frac{d-1}{2}}  \int_{-1}^{1} ( 1+ a x)^{v - \frac{d-1}{2}} ( 1- x^2)^{ \frac{d - 3}{2}}\d x \\
		& \leq c_d e^{v  t-\frac{d-1}{2}t} \left( I_+ + I_-\right),
	\end{align*}
	where 
	$$
	I_{\pm} = \int_{0}^{1} ( 1 \pm a x)^{v - \frac{d-1}{2}} ( 1- x^2)^{ \frac{d - 3}{2}}\d x.
	$$
	We have $I_+ \leq  \int_{0}^{1}   ( 1- x^2)^{ \frac{d - 3}{2}}\d x < \infty$. On the other hand, setting $a =  1- \delta$ and $b = a / \delta = (1 -\delta) / \delta$, we get through the change of variable $x \to 1 -x$:
	\begin{align*}
		I_- & = 	\int_{0}^{1} ( 1 - a (1-x) )^{v - \frac{d-1}{2}} ( x (2 -x) )^{ \frac{d - 3}{2}}\d x	\\
		& =  \delta^{v - \frac{d-1}{2}} \int_{0}^{1} ( 1 + b x )^{v - \frac{d-1}{2}} ( x (2 -x) )^{ \frac{d - 3}{2}}\d x \\
		& \leq C \delta^{v - \frac{d-1}{2}} \int_{0}^{1} ( 1 + b x )^{v - \frac{d-1}{2}}  x  ^{ \frac{d - 3}{2}}\d x, 
	\end{align*}		
	where at the last line, we have used that for all $x \in [0,1]$, $(2 -x)^{ \frac{d - 3}{2}} \leq C$ with $C = 1$ for $d \in \{2,3\}$ and $C = 2^{(d-3)/2}$ for $d \geq 4$. With the change of variable $y = bx$ and recalling that $\delta b = a \in [\tanh t_0,1)$, we find: 
	$$
	I_- \leq C a^{-\frac{d-1}{2}} \delta^{v}  \int_{0}^{b} ( 1 + y )^{v - \frac{d-1}{2}}  y  ^{ \frac{d - 3}{2}}\d y.
	$$
	We have $b = b(t) = (e^{2t} - 1)/ 2 \in [b_0,e^{2t}]$ with $b_0 = b(t_0) >0$. We write
	\begin{align*}
		\delta^{v}  \int_{0}^{b} ( 1 + y )^{v - \frac{d-1}{2}}  y  ^{ \frac{d - 3}{2}}\d y & = \delta^{v}  \int_{0}^{b_0} \dots +  \delta^{v}  \int_{b_0}^{b} \dots.
	\end{align*}
	Note that in a neighborhood of $0$, the integrand is equivalent to $y  ^{ \frac{d - 3}{2}}$, hence 
	$$
	\delta^{v}  \int_{0}^{b_0} ( 1 + y )^{v - \frac{d-1}{2}}  y  ^{ \frac{d - 3}{2}}\d y \leq C
	$$
	for some constant $C = C(d,t_0)$ (recall that $v \leq (d-1)/2$). Similarly, 
	$$
	\delta^{v}  \int_{b_0}^{b} ( 1 + y )^{v - \frac{d-1}{2}}  y  ^{ \frac{d - 3}{2}}\d y \leq \delta^{v}  \int_{b_0}^{b} ( 1 + y )^{v} y  ^{ -1} \d y \leq C \delta^{v}  \int_{b_0}^{b} y  ^{v -1} \d y,
	$$
	for some constant $C = C(d,t_0)$.  We finally use that
	$$
	\int_{b_0}^{b} y  ^{v -1} \d y  = \frac{1}{v}(b^v - b_0^v)= \frac{b^v}{v}\left(1 - \frac{b_0^v}{b^v}\right) \leq b^v \min \left( \frac{1}{v} , \ln \frac{b}{b_0}\right),
	$$
	Using again $\delta b = a \leq 1$ and $\ln (b) \sim 2t$ as $t \to \infty$, this proves \eqref{eq:bdnut1}.
	
	It remains to prove that if $u \in \R_+$, we have
	\begin{equation}\label{eq:bdnut2}
		|\nu_t(\lambda)|  \leq   \frac{C}{u^{\frac{d-1}{2}}}   e^{-\frac{d-1}{2}t}. 
	\end{equation}
	This bound improves  \eqref{eq:bdnut1} for large $u$. It is obtained by taking into account the oscillatory part in the formula for $\nu_t(\lambda)$. With $ a = 1- \delta = \tanh (t)$ as above, from Proposition \ref{prop:SpAt0} we write
	\begin{align*}
		\nu_t (\lambda)  &= c_d (\cosh t )^{iu - \frac{d-1}{2}}   \int_{-1}^{1} ( 1+ a x)^{iu - \frac{d-1}{2}} ( 1- x^2)^{ \frac{d - 3}{2}} \d x.
	\end{align*}

For $d \geq 4$, we start by performing integration by parts to take advantage of the factor $( 1- x^2)^{ \frac{d - 3}{2}}$. 
To this end, we fix $\beta \in \mathbb{C}$ and $\alpha \in [0,1)$. For integers $ n,k \geq 0$ with $2k \leq n$, let $\mathcal P_{k,n} \subset \mathbb{C}[X]$ be the set of polynomials of degree at most $n$ which are divisible by $(1-x^2)^k$.  For $P \in \mathbb{C}[X]$, we define
$$
I_k (P) = \int_{-1}^{1} ( 1+ a x)^{\beta - k-1} P(x) ( 1- x^2)^{ -\alpha} dx.
$$
By integration by parts, if $k \geq 1$ and $P \in \mathcal P_{k,n}$, we observe that 
\begin{align*}
I_k (P) & = \frac{1} {a (\beta -k)}   \int_{-1}^{1} ( 1+ a x)^{\beta - k}  \left (  -2 \alpha \frac{ x P(x)}{1-x^2} - P'(x) \right) ( 1- x^2)^{ -\alpha} dx \\
& = \frac{1} {a (\beta -k)} I_{k-1} (Q_1),
\end{align*}
	  where $Q_1 \in \mathcal P_{k-1,n-1}$. In particular, by iteration, we deduce that there exists a polynomial $Q_k\in \mathcal P_{0,n-k}$ such that 
$$
I_k (P) =  \frac{ I_0(Q_k)} { a^k  \prod_{j=0}^{k-1}  (\beta - k  + j) } . 
$$
Note that $Q_k$ does not depend on the parameters $(a,\beta)$.
	  We apply this observation to $k = \lceil (d-3)/2 \rceil$, $P(x) = (1-x^2)^{   k  }$, $\alpha = 0$ for odd $d \geq 5$, $\alpha = 1/2$ for even $d \geq 4$ and  $\beta = iu + \alpha$ (so that $\frac{d-3}{2} = k - \alpha$, $\frac{d-1}{2} = k - \alpha+1 $). We deduce that for some polynomial $Q$ of degree  at most $k$ we have:
	  				\begin{align*}
	  \nu_t (\lambda)  &= \frac{ c_d (\cosh t )^{iu - \frac{d-1}{2}}}  { a^k  \prod_{j=0}^{k-1}  (i u - \frac{d-3}{2} +j ) }  \int_{-1}^{1} ( 1+ a x)^{iu - 1 + \alpha} Q(x) (1-x^2)^{-\alpha}dx \\
	  & =  \frac{ c_d (\cosh t )^{iu - \frac{d-1}{2}}}  { a^k  \prod_{j=0}^{k-1}  (i u -  \frac{d-3}{2} +j ) } \left( J_+ + J_-\right),  
	   \end{align*}
		with
		$$
		J_{\pm} = \int_{0}^{1} ( 1 \pm a x)^{iu - 1 + \alpha} Q(x) ( 1- x^2)^{ -\alpha}dx.
		$$	
		The proof of \eqref{eq:bdnut2} will be complete if we prove that 
\begin{equation}\label{eq:bdnut2b}
		|J_{\pm}|  \leq   \frac{C}{u^{1-\alpha}}. 
\end{equation}
	We start with $J_-$ and set $b = a /  \delta$ as above. Since $((1-x)^l)_{l \geq 0}$ is a basis of $\C[X]$, by linearity, we may assume that $Q(x) = (1-x)^l$ for some integer $0 \leq l \leq k$. From what precedes, we get
	\begin{align*}
		J_- &=  \delta^{iu - 1 + \alpha} \int_{0}^{1} ( 1+ b x )^{iu - 1 + \alpha} x^l ( x(2- x))^{ -\alpha } \d x  \\
		& = \delta^{iu} u^{-1} (\delta b)^{-1 + \alpha } b ^{-l} \int_{0}^{u \ln (1 + b)} e^{iy } e^{y \alpha / u }  (e^{y/u} -1)^{l - \alpha}  (2-(e^{y/u} -1)/b )^{- \alpha } \d y,
	\end{align*}
	where we have performed the change of variable $y = u \ln ( 1 + bx)$, that is $x = (e^{y/u} -1)/b$. Since $\delta b = a$, we find
	\begin{align*}
		J_- 
		& = \delta^{iu} u^{-1+\alpha} a^{-1 + \alpha } b ^{-l} \int_{0}^{u \ln (1 + b)} e^{iy} y^{-\alpha} f_b\left(\frac y u \right) \d y,
	\end{align*}
	where we have set 
	$$
	f_b(x) =  \frac{(e^{x} -1)^{l} }{ \left( g(x)  h_b(x) \right)^{\alpha} }, \quad g(x) =  \frac{1-e^{-x} }{x} , \quad h_b(x) =  2 - \frac{e^{x} -1}{b}.
	$$		
	The functions $g,h_b$ and $x \mapsto (e^x-1)^l$ are analytic on $\C$. If $\alpha = 0$, we deduce that $f_b(x)$ is analytic on $\C$. If $\alpha = 1/2$, consider the principal branch of the function $z \mapsto z^{-\alpha}$ analytic on $\C \backslash \R^-$ with $\R^- = (-\infty,0]$.  For $\theta \geq 0$, let $C_\theta = \{ x \in \C : \mathrm{arg} (x) \in [-\theta,\theta] \}$ where $\mathrm{arg}(x) \in [-\pi,\pi)$ is the argument of a complex number $x$. Let $C^*_\theta = \C_\theta \backslash \{0\}$ and $C_\theta (r) = C_\theta \cap B(r)$ where $B(r) = \{ x : |x| \leq r\}$. We claim that the image of $C_{\pi/8}  (\ln (1+b))$ by $x \mapsto g(x)  h_b(x)$ is contained in $C^*_{7\pi / 8}$. This claim implies that  $f_b$ is analytic on $C_{\pi/8}  (\ln (1+b))$.

Let us prove the claim. First, for all $x \in C_{\pi/2} ( \ln(1+b))$, it is easy to check that  $|e^x - 1|/b \leq (e^{|x|} - 1 ) / b \leq 1$. Thus, for all $x \in C_{\pi/2} ( \ln(1+b))$, $h_b(x)$ is in a ball of radius $1$ centered at $2$. In particular, we have $h_b(x) \in C^*_{\pi/4}$. Note also that $g(0) = 1$ and $g(x) \ne 0$ for all $x \in \C$. Also, for $ x \in C^*_\theta$ with $\theta \leq \pi/2$, we have $|\mathrm{arg} (g(x))| \leq |\mathrm{arg} (1 - e^{-x})| + |\mathrm{arg} (x^{-1})| \leq \pi / 2 + \theta$ because $\Re(1 - e^{-x}) = 1 -e^{-|x| \cos (\theta)} \cos (|x| \sin (\theta))  >0$. So we have checked that for all $x \in C^*_\theta$, we have $g(x)  h_b(x) \in C^*_{\pi/2 + \pi/4 + \theta}$. Since $g(0)  h_b(0) = 2$, the latter is true for all $x \in C_\theta$. We deduce the claim. 


We can now complete the proof of \eqref{eq:bdnut2b} for $J_-$. Since  $f_b$ is analytic on $C_{\pi/8}  (\ln (1+b))$, we can apply Cauchy formula around a counterclockwise contour around the set $[0,u \ln (1+b)] \cup \{ x : x =  u \ln (1+b) e^{i \theta} , \theta \in [0,\pi/8]\} \cup e^{i \pi/8} [0,u \ln (1+b)]$. We get, with $r = u \ln (1+ b)$,
\begin{align*}
	K_- & = \int_{0}^{r} e^{iy} y^{-\alpha} f_b\left(\frac y u \right) \d y \\
	& = \int_{0}^{r} e^{iy  e^{i\pi/8}} ( e^{i\pi/8} y)^{-\alpha} f_b\left(\frac {y  e^{i\pi/8}}{ u}  \right) e^{i\pi/8}\d y -  \int_{0}^{\pi/8} e^{iy e^{i \theta} } (r e^{i\theta}) ^{-\alpha} f_b\left(\frac {r e^{i \theta}}{ u}  \right) r \d \theta.
\end{align*}
Taking absolute values, with $c = \sin (\pi/8) > 0$, we get 
$$
| K_- | \leq \int_{0}^{r} e^{-c y } y^{-\alpha} \left| f_b\left(\frac {y  e^{i\pi/8}}{ u}  \right)\right|  \d y + r^{1-\alpha} \int_{0}^{\pi/8} e^{- r \sin ( \theta) }  \left|f_b\left(\frac {r e^{i \theta}}{ u}  \right)\right|  \d \theta.
$$
Next we can use $\sin(\theta ) \geq 2 \theta / \pi$ and that for all $x \in C_{\pi/2}( \ln(1+b))$ we have $|f_b(x)| \leq C b^l$  for some constant $C > 0$. We get that
$$
| K_- | \leq C  b^l \int_0^\infty e^{-cy} y^{-\alpha} \d y + C b^l r^{1-\alpha}  \int_{0}^{\pi/8} e^{- 2 r \theta / \pi } \d\theta \leq C' b^l ( 1+ r^{-\alpha}),  
$$
for some new constant $C'$. Since $J_- = \delta^{iu} u^{-1+\alpha} a^{-1 + \alpha } b ^{-l} K_-$, this completes the proof of \eqref{eq:bdnut2b} for $J_-$. 


It remains to prove \eqref{eq:bdnut2} for $J_+$. The proof is the same as the proof for $J_-$ so we will use the same notation. Using the same argument as above, we assume that $Q(x) = (1-x)^l$. We now set  $\delta = (1+a)$ and $b = a / (1 + a) \in [b_0,1/2)$ with $b _0 = b(t_0) \in (0,1/2)$. We write 
\begin{align*}
	J_+ &=  \delta^{iu - 1 + \alpha} \int_{0}^{1} ( 1 - b x )^{iu - 1 + \alpha} x^l ( x(2- x))^{ -\alpha } \d x  \\
	& = \delta^{iu} u^{-1} (\delta b)^{-1 + \alpha } b ^{-l} \int_{0}^{- u \ln (1 - b)} e^{-iy } e^{-y \alpha / u }  (1-e^{-y/u})^{l - \alpha}  (2-(1 - e^{-y/u})/b )^{- \alpha } \d y,
\end{align*}
where we have performed the change of variable $y =  - u \ln ( 1 - bx)$, that is $x = (1 - e^{-y/u})/b$. Since $\delta b = a$, we find
\begin{align*}
	J_+ 
	& = \delta^{iu} u^{-1+\alpha} a^{-1 + \alpha } b ^{-l} \int_{0}^{- u \ln (1 - b)} e^{-iy} y^{-\alpha} f_b\left(\frac y u \right) \d y,
\end{align*}
where now we have
$$
f_b(x) =  \frac{(1-e^{-x})^{l} }{ \left( g(x)  h_b(x) \right)^{\alpha} }, \quad g(x) =  \frac{e^{x} -1 }{x} , \quad h_b(x) =  2 - \frac{1- e^{-x}}{b}.
$$		
We can then conclude similarly thanks to a contour integration around the set $[0, - u \ln (1 - b)] \cup \{ x : x =  -u \ln (1+b) e^{ - i \theta} , \theta \in [0,\pi/8]\} \cup e^{ - i \pi/8} [0, - u \ln (1-b)]$ (note that in this case, it is easier since $\ln (1-b) $ and $b$ are uniformly bounded by some $C = C(t_0)$). 
\end{proof}

\end{document}
